% =================================================================
% TESTING PLAN - MODIFICA DI UNA ISSUE
% =================================================================

\subsection{Testing Plan: Modifica di una Issue}
\label{sec:testing-modifica-issue}

Di seguito sono riportati i test case relativi alla funzionalit\`a di modifica di una issue.
Per ogni scenario vengono presentate due tabelle: una per il frontend e una per il backend.

% -----------------------------------------------------------------
% SCENARIO 1: Modifica dello stato dell'issue
% -----------------------------------------------------------------
\subsubsection{Modifica dello stato}

\paragraph{Frontend}

{
\small
\begin{longtable}{|c|p{3.8cm}|p{3.8cm}|p{3.8cm}|p{3.8cm}|}
\hline
\textbf{ID} & \textbf{Precondizioni} & \textbf{Input} & \textbf{Output} & \textbf{Postcondizioni} \\
\hline
\endfirsthead
\hline
\textbf{ID} & \textbf{Precondizioni} & \textbf{Input} & \textbf{Output} & \textbf{Postcondizioni} \\
\hline
\endhead
1 &
L'utente \`e autenticato come amministratore. Il modale dell'issue \`e aperto. &
L'utente seleziona un nuovo stato dalla dropdown (es. ``In Progress''). &
Viene mostrato un toast di successo con il messaggio ``Stato cambiato in In Progress''. Il modale si chiude. &
Lo stato dell'issue \`e aggiornato nella lista delle issue. \\
\hline
2 &
L'utente \`e autenticato come assegnatario dell'issue. L'issue ha stato ``To Do''. &
L'utente seleziona ``In Progress'' dalla dropdown dello stato. &
Viene mostrato un toast di successo. Il modale si chiude. &
Lo stato dell'issue \`e aggiornato nella lista. \\
\hline
4 &
L'utente \`e autenticato ma non \`e n\'e amministratore n\'e assegnatario n\'e autore dell'issue. &
L'utente tenta di cambiare lo stato. &
La dropdown dello stato \`e disabilitata. L'utente non pu\`o interagire con essa. &
Lo stato dell'issue resta invariato. \\
\hline
6 &
L'utente \`e autenticato come amministratore. &
L'utente seleziona un nuovo stato ma il server restituisce un errore. &
Viene mostrato un toast di errore con il messaggio ``Errore durante l'aggiornamento''. Lo stato nell'interfaccia torna al valore precedente. &
Lo stato dell'issue resta invariato. \\
\hline
\end{longtable}
}
\newpage
\paragraph{Backend}

{
\small
\begin{longtable}{|c|p{3.8cm}|p{3.8cm}|p{3.8cm}|p{3.8cm}|}
\hline
\textbf{ID} & \textbf{Precondizioni} & \textbf{Input} & \textbf{Output} & \textbf{Postcondizioni} \\
\hline
\endfirsthead
\hline
\textbf{ID} & \textbf{Precondizioni} & \textbf{Input} & \textbf{Output} & \textbf{Postcondizioni} \\
\hline
\endhead
7 &
L'issue con id 1 esiste nel database. L'utente richiedente \`e l'amministratore. Viene effettuata una richiesta PUT a \texttt{/api/issues/1}. &
Body: \texttt{\{``statoIssue'': ``IN\_LAVORAZIONE''\}}. &
Risposta HTTP 200 con il DTO dell'issue aggiornata contenente \texttt{statoIssue = IN\_LAVORAZIONE}. &
Nel database il campo \texttt{stato\_issue} dell'issue 1 \`e aggiornato a \texttt{IN\_LAVORAZIONE}. La data di ultimo aggiornamento \`e aggiornata. \\
\hline
8 &
L'issue con id 1 esiste nel database. L'utente richiedente \`e l'assegnatario. Viene effettuata una richiesta PUT a \texttt{/api/issues/1}. &
Body: \texttt{\{``statoIssue'': ``IN\_REVISIONE''\}}. &
Risposta HTTP 200 con il DTO dell'issue aggiornata. &
Lo stato nel database \`e aggiornato a \texttt{IN\_REVISIONE}. \\
\hline
9 &
L'issue con id 1 esiste nel database. L'utente richiedente \`e l'autore dell'issue. Viene effettuata una richiesta PUT a \texttt{/api/issues/1}. &
Body: \texttt{\{``statoIssue'': ``TODO''\}}. &
Risposta HTTP 200 con il DTO dell'issue aggiornata. &
Lo stato nel database \`e aggiornato a \texttt{TODO}. \\
\hline
10 &
L'issue con id 1 esiste nel database. L'utente richiedente non \`e n\'e amministratore, n\'e assegnatario, n\'e autore. Viene effettuata una richiesta PUT a \texttt{/api/issues/1}. &
Body: \texttt{\{``statoIssue'': ``RISOLTA''\}}. &
Risposta HTTP 403 con messaggio ``Non hai i permessi per modificare questa issue.''. &
Lo stato dell'issue nel database resta invariato. \\
\hline
11 &
L'issue con id 999 non esiste nel database. Viene effettuata una richiesta PUT a \texttt{/api/issues/999}. &
Body: \texttt{\{``statoIssue'': ``TODO''\}}. &
Risposta HTTP 404 con messaggio ``Issue non trovata''. &
Nessuna modifica nel database. \\
\hline
\end{longtable}
}

% -----------------------------------------------------------------
% SCENARIO 2: Modifica dell'assegnatario
% -----------------------------------------------------------------
\subsubsection{Modifica dell'assegnatario}

\paragraph{Frontend}

{
\small
\begin{longtable}{|c|p{3.8cm}|p{3.8cm}|p{3.8cm}|p{3.8cm}|}
\hline
\textbf{ID} & \textbf{Precondizioni} & \textbf{Input} & \textbf{Output} & \textbf{Postcondizioni} \\
\hline
\endfirsthead
\hline
\textbf{ID} & \textbf{Precondizioni} & \textbf{Input} & \textbf{Output} & \textbf{Postcondizioni} \\
\hline
\endhead
12 &
L'utente \`e autenticato come amministratore. Il modale dell'issue \`e aperto. La lista dei membri del progetto \`e stata caricata. &
L'utente seleziona un membro del progetto dalla dropdown dell'assegnatario. &
Viene mostrato un toast di successo con il messaggio ``Issue assegnata a [nome]''. &
L'assegnatario dell'issue \`e aggiornato nell'interfaccia. \\
\hline
13 &
L'utente \`e autenticato come amministratore. L'issue ha gi\`a un assegnatario. &
L'utente seleziona ``Non assegnato'' dalla dropdown dell'assegnatario. &
Viene mostrato un toast di successo con il messaggio ``Issue non assegnata''. &
L'assegnatario dell'issue viene rimosso nell'interfaccia. \\
\hline
14 &
L'utente \`e autenticato ma non \`e amministratore. &
L'utente visualizza la dropdown dell'assegnatario. &
La dropdown \`e disabilitata. L'utente non pu\`o interagire con essa. &
L'assegnatario resta invariato. \\
\hline
15 &
L'utente \`e autenticato come amministratore. &
L'utente seleziona un assegnatario ma il server restituisce un errore. &
Viene mostrato un toast di errore con il messaggio ``Errore durante l'aggiornamento''. L'assegnatario nell'interfaccia torna al valore precedente. &
L'assegnatario dell'issue resta invariato. \\
\hline
\end{longtable}
}

\newpage
\paragraph{Backend}
{
\small
\begin{longtable}{|c|p{3.8cm}|p{3.8cm}|p{3.8cm}|p{3.8cm}|}
\hline
\textbf{ID} & \textbf{Precondizioni} & \textbf{Input} & \textbf{Output} & \textbf{Postcondizioni} \\
\hline
\endfirsthead
\hline
\textbf{ID} & \textbf{Precondizioni} & \textbf{Input} & \textbf{Output} & \textbf{Postcondizioni} \\
\hline
\endhead
16 &
L'issue con id 1 esiste. L'utente richiedente \`e amministratore. L'utente ``mario@test.it'' \`e membro del progetto. Viene effettuata una richiesta PUT a \texttt{/api/issues/1}. &
Body: \texttt{\{``assegnatario'': ``mario@test.it''\}}. &
Risposta HTTP 200 con il DTO dell'issue aggiornata con l'assegnatario impostato. &
Nel database l'assegnatario dell'issue 1 \`e ``mario@test.it''. Viene creata una notifica di assegnazione per l'utente. \\
\hline
17 &
L'issue con id 1 esiste e ha gi\`a un assegnatario. L'utente richiedente \`e amministratore. Viene effettuata una richiesta PUT a \texttt{/api/issues/1}. &
Body: \texttt{\{``assegnatario'': ``''\}}. &
Risposta HTTP 200 con il DTO dell'issue aggiornata con assegnatario nullo. &
Nel database l'assegnatario dell'issue 1 \`e impostato a NULL. \\
\hline
18 &
L'issue con id 1 esiste. L'utente richiedente \`e amministratore. L'utente ``esterno@test.it'' non \`e membro del progetto e non \`e amministratore. Viene effettuata una richiesta PUT a \texttt{/api/issues/1}. &
Body: \texttt{\{``assegnatario'': ``esterno@test.it''\}}. &
Risposta HTTP 400 con messaggio ``L'utente selezionato non \`e membro del progetto''. &
L'assegnatario dell'issue nel database resta invariato. \\
\hline
19 &
L'issue con id 1 esiste. L'utente richiedente \`e amministratore. Viene effettuata una richiesta PUT a \texttt{/api/issues/1}. &
Body: \texttt{\{``assegnatario'': ``inesistente@test.it''\}}. &
Risposta HTTP 404 con messaggio ``Assegnatario non trovato''. &
L'assegnatario dell'issue nel database resta invariato. \\
\hline
20 &
L'issue con id 1 esiste. L'utente richiedente \`e amministratore. L'issue ha gi\`a assegnatario ``mario@test.it''. Viene effettuata una richiesta PUT a \texttt{/api/issues/1}. &
Body: \texttt{\{``assegnatario'': ``mario@test.it''\}}. &
Risposta HTTP 200 con il DTO dell'issue. &
L'assegnatario resta ``mario@test.it''. Non viene creata una nuova notifica perch\'e l'assegnatario non \`e cambiato. \\
\hline
21 &
L'issue con id 1 esiste. L'utente richiedente non \`e n\'e amministratore, n\'e assegnatario, n\'e autore dell'issue. Viene effettuata una richiesta PUT a \texttt{/api/issues/1}. &
Body: \texttt{\{``assegnatario'': ``mario@test.it''\}}. &
Risposta HTTP 403 con messaggio ``Non hai i permessi per modificare questa issue.''. &
L'assegnatario nel database resta invariato. \\
\hline
\end{longtable}
}

% -----------------------------------------------------------------
% SCENARIO 3: Modifica della priorit\`a
% -----------------------------------------------------------------
\subsubsection{Modifica della priorit\`a}

\paragraph{Frontend}

{
\small
\begin{longtable}{|c|p{3.8cm}|p{3.8cm}|p{3.8cm}|p{3.8cm}|}
\hline
\textbf{ID} & \textbf{Precondizioni} & \textbf{Input} & \textbf{Output} & \textbf{Postcondizioni} \\
\hline
\endfirsthead
\hline
\textbf{ID} & \textbf{Precondizioni} & \textbf{Input} & \textbf{Output} & \textbf{Postcondizioni} \\
\hline
\endhead
22 &
L'utente \`e autenticato ed \`e autorizzato a modificare l'issue (amministratore, assegnatario o autore). Il modale dell'issue \`e aperto. &
L'utente seleziona ``Alta'' dalla dropdown della priorit\`a. &
Viene mostrato un toast di successo con il messaggio ``Priorit\`a cambiata in Alta''. &
La priorit\`a dell'issue \`e aggiornata nell'interfaccia. \\
\hline
23 &
L'utente \`e autenticato ma non ha i permessi di modifica sull'issue. &
L'utente visualizza la dropdown della priorit\`a. &
La dropdown \`e disabilitata. L'utente non pu\`o selezionare nulla. &
La priorit\`a resta invariata. \\
\hline
24 &
L'utente \`e autenticato e autorizzato. &
L'utente cambia la priorit\`a ma il server restituisce un errore. &
Viene mostrato un toast di errore con il messaggio ``Errore durante l'aggiornamento''. La priorit\`a nell'interfaccia torna al valore precedente. &
La priorit\`a dell'issue resta invariata. \\
\hline
\end{longtable}
}

\paragraph{Backend}

{
\small
\begin{longtable}{|c|p{3.8cm}|p{3.8cm}|p{3.8cm}|p{3.8cm}|}
\hline
\textbf{ID} & \textbf{Precondizioni} & \textbf{Input} & \textbf{Output} & \textbf{Postcondizioni} \\
\hline
\endfirsthead
\hline
\textbf{ID} & \textbf{Precondizioni} & \textbf{Input} & \textbf{Output} & \textbf{Postcondizioni} \\
\hline
\endhead
25 &
L'issue con id 1 esiste. L'utente richiedente \`e amministratore. Viene effettuata una richiesta PUT a \texttt{/api/issues/1}. &
Body: \texttt{\{``prioritaIssue'': ``ALTA''\}}. &
Risposta HTTP 200 con il DTO dell'issue aggiornata con \texttt{prioritaIssue = ALTA}. &
Nel database la priorit\`a dell'issue 1 \`e \texttt{ALTA}. \\
\hline
26 &
L'issue con id 1 esiste. L'utente richiedente \`e l'assegnatario. Viene effettuata una richiesta PUT a \texttt{/api/issues/1}. &
Body: \texttt{\{``prioritaIssue'': ``CRITICA''\}}. &
Risposta HTTP 200 con il DTO aggiornato. &
La priorit\`a nel database \`e \texttt{CRITICA}. \\
\hline
27 &
L'issue con id 1 esiste. L'utente richiedente non ha i permessi. Viene effettuata una richiesta PUT a \texttt{/api/issues/1}. &
Body: \texttt{\{``prioritaIssue'': ``MEDIA''\}}. &
Risposta HTTP 403 con messaggio ``Non hai i permessi per modificare questa issue.''. &
La priorit\`a nel database resta invariata. \\
\hline
\end{longtable}
}

% -----------------------------------------------------------------
% SCENARIO 4: Modifica delle etichette
% -----------------------------------------------------------------
\subsubsection{Modifica delle etichette}

\paragraph{Frontend}

{
\small
\begin{longtable}{|c|p{3.8cm}|p{3.8cm}|p{3.8cm}|p{3.8cm}|}
\hline
\textbf{ID} & \textbf{Precondizioni} & \textbf{Input} & \textbf{Output} & \textbf{Postcondizioni} \\
\hline
\endfirsthead
\hline
\textbf{ID} & \textbf{Precondizioni} & \textbf{Input} & \textbf{Output} & \textbf{Postcondizioni} \\
\hline
\endhead
28 &
L'utente \`e autenticato e autorizzato. Il modale dell'issue \`e aperto. &
L'utente digita una nuova etichetta nel campo di input e preme Invio. &
L'etichetta compare nella lista delle etichette. Dopo 800ms viene mostrato un toast di successo ``Etichette aggiornate''. &
Le etichette dell'issue sono aggiornate nell'interfaccia e nel backend. \\
\hline
29 &
L'utente \`e autenticato e autorizzato. L'issue ha gi\`a delle etichette. &
L'utente rimuove un'etichetta esistente cliccando sulla X dell'etichetta. &
L'etichetta viene rimossa dalla lista. Dopo 800ms viene mostrato un toast di successo ``Etichette aggiornate''. &
Le etichette dell'issue sono aggiornate. \\
\hline
30 &
L'utente \`e autenticato ma non ha i permessi di modifica sull'issue. &
L'utente visualizza la sezione etichette. &
Le etichette vengono mostrate in sola lettura, senza il campo di input. &
Nessuna modifica possibile. \\
\hline
31 &
L'utente \`e autenticato e autorizzato. L'issue non ha etichette. &
L'utente visualizza la sezione etichette. &
Viene mostrato il messaggio ``Nessuna etichetta'' insieme al campo di input per aggiungerne. &
Nessuna modifica fino a che l'utente non aggiunge un'etichetta. \\
\hline
32 &
L'utente \`e autenticato e autorizzato. &
L'utente aggiunge un'etichetta ma il server restituisce un errore. &
Viene mostrato un toast di errore con il messaggio ``Errore durante l'aggiornamento''. Le etichette nell'interfaccia tornano al valore precedente. &
Le etichette dell'issue restano invariate. \\
\hline
\end{longtable}
}

\paragraph{Backend}

{
\small
\begin{longtable}{|c|p{3.8cm}|p{3.8cm}|p{3.8cm}|p{3.8cm}|}
\hline
\textbf{ID} & \textbf{Precondizioni} & \textbf{Input} & \textbf{Output} & \textbf{Postcondizioni} \\
\hline
\endfirsthead
\hline
\textbf{ID} & \textbf{Precondizioni} & \textbf{Input} & \textbf{Output} & \textbf{Postcondizioni} \\
\hline
\endhead
33 &
L'issue con id 1 esiste. L'utente richiedente \`e amministratore. Viene effettuata una richiesta PUT a \texttt{/api/issues/1}. &
Body: \texttt{\{``etichette'': [``bugfix'', ``urgente'']\}}. &
Risposta HTTP 200 con il DTO dell'issue aggiornata. &
Nel database le etichette ``bugfix'' e ``urgente'' sono associate all'issue 1. Se non esistevano, vengono create. \\
\hline
34 &
L'issue con id 1 esiste. L'utente richiedente \`e l'autore. Viene effettuata una richiesta PUT a \texttt{/api/issues/1}. &
Body: \texttt{\{``etichette'': [``nuova-etichetta'']\}}. &
Risposta HTTP 200 con il DTO aggiornato. &
L'etichetta ``nuova-etichetta'' viene creata nel progetto (se non esiste) e associata all'issue. \\
\hline
35 &
L'issue con id 1 esiste. L'utente richiedente non ha i permessi. Viene effettuata una richiesta PUT a \texttt{/api/issues/1}. &
Body: \texttt{\{``etichette'': [``test'']\}}. &
Risposta HTTP 403 con messaggio ``Non hai i permessi per modificare questa issue.''. &
Nessuna etichetta viene aggiunta nel database. \\
\hline
\end{longtable}
}

% -----------------------------------------------------------------
% SCENARIO 5: Modifica del titolo
% -----------------------------------------------------------------
\subsubsection{Modifica del titolo}

\paragraph{Frontend}

{
\small
\begin{longtable}{|c|p{3.8cm}|p{3.8cm}|p{3.8cm}|p{3.8cm}|}
\hline
\textbf{ID} & \textbf{Precondizioni} & \textbf{Input} & \textbf{Output} & \textbf{Postcondizioni} \\
\hline
\endfirsthead
\hline
\textbf{ID} & \textbf{Precondizioni} & \textbf{Input} & \textbf{Output} & \textbf{Postcondizioni} \\
\hline
\endhead
36 &
L'utente \`e autenticato e autorizzato. Il modale dell'issue \`e aperto. &
L'utente modifica il titolo dell'issue e conferma. &
Viene mostrato un toast di successo. Il titolo viene aggiornato nel modale e nella lista. &
Il titolo dell'issue \`e aggiornato nell'interfaccia. \\
\hline
37 &
L'utente \`e autenticato e autorizzato. &
L'utente modifica il titolo ma il server restituisce un errore. &
Viene mostrato un toast di errore con il messaggio ``Errore durante l'aggiornamento''. Il titolo torna al valore precedente. &
Il titolo dell'issue resta invariato. \\
\hline
\end{longtable}
}

\paragraph{Backend}

{
\small
\begin{longtable}{|c|p{3.8cm}|p{3.8cm}|p{3.8cm}|p{3.8cm}|}
\hline
\textbf{ID} & \textbf{Precondizioni} & \textbf{Input} & \textbf{Output} & \textbf{Postcondizioni} \\
\hline
\endfirsthead
\hline
\textbf{ID} & \textbf{Precondizioni} & \textbf{Input} & \textbf{Output} & \textbf{Postcondizioni} \\
\hline
\endhead
38 &
L'issue con id 1 esiste. L'utente richiedente \`e amministratore. Viene effettuata una richiesta PUT a \texttt{/api/issues/1}. &
Body: \texttt{\{``titolo'': ``Nuovo titolo''\}}. &
Risposta HTTP 200 con il DTO dell'issue con il titolo aggiornato. &
Nel database il titolo dell'issue 1 \`e ``Nuovo titolo''. La data di ultimo aggiornamento \`e aggiornata. \\
\hline
39 &
L'issue con id 1 esiste. L'utente richiedente \`e l'autore. Viene effettuata una richiesta PUT a \texttt{/api/issues/1}. &
Body: \texttt{\{``titolo'': ``Titolo corretto''\}}. &
Risposta HTTP 200 con il DTO aggiornato. &
Il titolo nel database \`e ``Titolo corretto''. \\
\hline
40 &
L'issue con id 1 esiste. L'utente richiedente non ha i permessi. Viene effettuata una richiesta PUT a \texttt{/api/issues/1}. &
Body: \texttt{\{``titolo'': ``Test''\}}. &
Risposta HTTP 403 con messaggio ``Non hai i permessi per modificare questa issue.''. &
Il titolo nel database resta invariato. \\
\hline
\end{longtable}
}

% -----------------------------------------------------------------
% SCENARIO 6: Modifica della descrizione
% -----------------------------------------------------------------
\subsubsection{Modifica della descrizione}

\paragraph{Frontend}

{
\small
\begin{longtable}{|c|p{3.8cm}|p{3.8cm}|p{3.8cm}|p{3.8cm}|}
\hline
\textbf{ID} & \textbf{Precondizioni} & \textbf{Input} & \textbf{Output} & \textbf{Postcondizioni} \\
\hline
\endfirsthead
\hline
\textbf{ID} & \textbf{Precondizioni} & \textbf{Input} & \textbf{Output} & \textbf{Postcondizioni} \\
\hline
\endhead
41 &
L'utente \`e autenticato e autorizzato. Il modale dell'issue \`e aperto. &
L'utente modifica la descrizione dell'issue e conferma. &
Viene mostrato un toast di successo. La descrizione viene aggiornata nel modale. &
La descrizione dell'issue \`e aggiornata nell'interfaccia. \\
\hline
42 &
L'utente \`e autenticato e autorizzato. &
L'utente modifica la descrizione ma il server restituisce un errore. &
Viene mostrato un toast di errore. La descrizione torna al valore precedente. &
La descrizione dell'issue resta invariata. \\
\hline
\end{longtable}
}

\paragraph{Backend}

{
\small
\begin{longtable}{|c|p{3.8cm}|p{3.8cm}|p{3.8cm}|p{3.8cm}|}
\hline
\textbf{ID} & \textbf{Precondizioni} & \textbf{Input} & \textbf{Output} & \textbf{Postcondizioni} \\
\hline
\endfirsthead
\hline
\textbf{ID} & \textbf{Precondizioni} & \textbf{Input} & \textbf{Output} & \textbf{Postcondizioni} \\
\hline
\endhead
43 &
L'issue con id 1 esiste. L'utente richiedente \`e amministratore. Viene effettuata una richiesta PUT a \texttt{/api/issues/1}. &
Body: \texttt{\{``descrizione'': ``Nuova descrizione dettagliata''\}}. &
Risposta HTTP 200 con il DTO dell'issue con la descrizione aggiornata. &
Nel database la descrizione dell'issue 1 \`e aggiornata. La data di ultimo aggiornamento \`e aggiornata. \\
\hline
44 &
L'issue con id 1 esiste. L'utente richiedente non ha i permessi. Viene effettuata una richiesta PUT a \texttt{/api/issues/1}. &
Body: \texttt{\{``descrizione'': ``Test''\}}. &
Risposta HTTP 403 con messaggio ``Non hai i permessi per modificare questa issue.''. &
La descrizione nel database resta invariata. \\
\hline
\end{longtable}
}

% -----------------------------------------------------------------
% SCENARIO 7: Modifica del tipo di issue
% -----------------------------------------------------------------
\subsubsection{Modifica del tipo di issue}

\paragraph{Frontend}

{
\small
\begin{longtable}{|c|p{3.8cm}|p{3.8cm}|p{3.8cm}|p{3.8cm}|}
\hline
\textbf{ID} & \textbf{Precondizioni} & \textbf{Input} & \textbf{Output} & \textbf{Postcondizioni} \\
\hline
\endfirsthead
\hline
\textbf{ID} & \textbf{Precondizioni} & \textbf{Input} & \textbf{Output} & \textbf{Postcondizioni} \\
\hline
\endhead
45 &
L'utente \`e autenticato e autorizzato. Il modale dell'issue \`e aperto. &
L'utente seleziona un nuovo tipo (es. ``Feature'') dalla dropdown del tipo. &
Viene mostrato un toast di successo. Il tipo viene aggiornato nell'interfaccia. &
Il tipo dell'issue \`e aggiornato. \\
\hline
46 &
L'utente \`e autenticato e autorizzato. &
L'utente cambia il tipo ma il server restituisce un errore. &
Viene mostrato un toast di errore. Il tipo torna al valore precedente. &
Il tipo dell'issue resta invariato. \\
\hline
\end{longtable}
}

\paragraph{Backend}

{
\small
\begin{longtable}{|c|p{3.8cm}|p{3.8cm}|p{3.8cm}|p{3.8cm}|}
\hline
\textbf{ID} & \textbf{Precondizioni} & \textbf{Input} & \textbf{Output} & \textbf{Postcondizioni} \\
\hline
\endfirsthead
\hline
\textbf{ID} & \textbf{Precondizioni} & \textbf{Input} & \textbf{Output} & \textbf{Postcondizioni} \\
\hline
\endhead
47 &
L'issue con id 1 esiste. L'utente richiedente \`e amministratore. Viene effettuata una richiesta PUT a \texttt{/api/issues/1}. &
Body: \texttt{\{``tipoIssue'': ``FEATURE''\}}. &
Risposta HTTP 200 con il DTO dell'issue con \texttt{tipoIssue = FEATURE}. &
Nel database il tipo dell'issue 1 \`e \texttt{FEATURE}. \\
\hline
48 &
L'issue con id 1 esiste. L'utente richiedente non ha i permessi. Viene effettuata una richiesta PUT a \texttt{/api/issues/1}. &
Body: \texttt{\{``tipoIssue'': ``BUG''\}}. &
Risposta HTTP 403 con messaggio ``Non hai i permessi per modificare questa issue.''. &
Il tipo nel database resta invariato. \\
\hline
\end{longtable}
}

% -----------------------------------------------------------------
% SCENARIO 8: Modifica di pi\`u campi contemporaneamente
% -----------------------------------------------------------------
\subsubsection{Modifica di pi\`u campi contemporaneamente}

\paragraph{Frontend}

{
\small
\begin{longtable}{|c|p{3.8cm}|p{3.8cm}|p{3.8cm}|p{3.8cm}|}
\hline
\textbf{ID} & \textbf{Precondizioni} & \textbf{Input} & \textbf{Output} & \textbf{Postcondizioni} \\
\hline
\endfirsthead
\hline
\textbf{ID} & \textbf{Precondizioni} & \textbf{Input} & \textbf{Output} & \textbf{Postcondizioni} \\
\hline
\endhead
49 &
L'utente \`e autenticato come amministratore. Il modale dell'issue \`e aperto. &
L'utente cambia prima la priorit\`a, poi l'assegnatario, poi lo stato, facendo una modifica alla volta. &
Per ogni modifica viene mostrato un toast di successo. Le modifiche sono visibili nell'interfaccia. &
L'issue nell'interfaccia riflette tutte le modifiche fatte. \\
\hline
\end{longtable}
}

\paragraph{Backend}

{
\small
\begin{longtable}{|c|p{3.8cm}|p{3.8cm}|p{3.8cm}|p{3.8cm}|}
\hline
\textbf{ID} & \textbf{Precondizioni} & \textbf{Input} & \textbf{Output} & \textbf{Postcondizioni} \\
\hline
\endfirsthead
\hline
\textbf{ID} & \textbf{Precondizioni} & \textbf{Input} & \textbf{Output} & \textbf{Postcondizioni} \\
\hline
\endhead
50 &
L'issue con id 1 esiste. L'utente richiedente \`e amministratore. Viene effettuata una richiesta PUT a \texttt{/api/issues/1}. &
Body: \texttt{\{``statoIssue'': ``IN\_LAVORAZIONE'', ``prioritaIssue'': ``ALTA'', ``assegnatario'': ``mario@test.it''\}}. &
Risposta HTTP 200 con il DTO dell'issue con tutti i campi aggiornati. &
Nel database stato, priorit\`a e assegnatario dell'issue 1 sono aggiornati. Viene creata una notifica di assegnazione. La data di ultimo aggiornamento \`e aggiornata. \\
\hline
\end{longtable}
}

% -----------------------------------------------------------------
% SCENARIO 9: Notifiche dopo la modifica
% -----------------------------------------------------------------
\subsubsection{Notifiche dopo la modifica}

\paragraph{Frontend}

{
\small
\begin{longtable}{|c|p{3.8cm}|p{3.8cm}|p{3.8cm}|p{3.8cm}|}
\hline
\textbf{ID} & \textbf{Precondizioni} & \textbf{Input} & \textbf{Output} & \textbf{Postcondizioni} \\
\hline
\endfirsthead
\hline
\textbf{ID} & \textbf{Precondizioni} & \textbf{Input} & \textbf{Output} & \textbf{Postcondizioni} \\
\hline
\endhead
51 &
L'utente A \`e autenticato come amministratore. L'utente B \`e membro del progetto ed \`e collegato all'applicazione. &
L'utente A modifica lo stato di un'issue del progetto. &
L'utente B riceve un aggiornamento in tempo reale tramite SSE. La lista delle issue dell'utente B si aggiorna automaticamente. &
L'utente B vede l'issue con lo stato aggiornato senza ricaricare la pagina. \\
\hline
\end{longtable}
}

\paragraph{Backend}

{
\small
\begin{longtable}{|c|p{3.8cm}|p{3.8cm}|p{3.8cm}|p{3.8cm}|}
\hline
\textbf{ID} & \textbf{Precondizioni} & \textbf{Input} & \textbf{Output} & \textbf{Postcondizioni} \\
\hline
\endfirsthead
\hline
\textbf{ID} & \textbf{Precondizioni} & \textbf{Input} & \textbf{Output} & \textbf{Postcondizioni} \\
\hline
\endhead
52 &
L'issue con id 1 esiste nel progetto P1. Il progetto P1 ha 3 membri. Viene effettuata una qualsiasi modifica all'issue 1. &
Qualsiasi campo modificato nel body della richiesta PUT. &
Il server invia un segnale SSE di tipo ``issue-update'' a tutti i membri del progetto P1 e a tutti gli amministratori. &
Tutti i membri del progetto e gli amministratori hanno ricevuto il segnale di aggiornamento. \\
\hline
53 &
L'issue con id 1 esiste. L'utente richiedente \`e amministratore. L'issue non aveva un assegnatario. Viene effettuata una richiesta PUT a \texttt{/api/issues/1}. &
Body: \texttt{\{``assegnatario'': ``mario@test.it''\}}. &
Risposta HTTP 200. Viene creata una notifica di assegnazione con testo ``Ti \`e stata assegnata l'issue: [titolo]'' per l'utente ``mario@test.it''. &
Nel database esiste una nuova notifica di assegnazione per ``mario@test.it''. \\
\hline
\end{longtable}
}
